\section{A logical repair for the weak-cover hybrid}\label{sec:pointed-repair}
\begingroup
\newcommand{\Sub}{\operatorname{Sub}}
\newcommand{\Eq}{\mathsf E}
\newcommand{\Mem}{\mathsf M}
\newcommand{\Low}{\mathcal L}
\newcommand{\Upp}{\mathcal U}

This section supplies a mathematical component for a possible hybrid of
U208's weak-cover emitter and PA196's equality-vacuity reduction. It does
\emph{not} supply a smaller transition table or an internal PA
consistency biconditional. The distinction between ordinary first-order
logic and the numeric machine calculus matters: the latter must establish
its quantifier and equality rules from its emitted heads.

We use the element-first membership notation $\Mem(x,s)$ of the U208
role-emitter audit: $x$ belongs to $s$. A coherent reversal in the numeric
encoding of the two arguments is a separate coding change. The primitive
symbol $\Eq$ below is not assumed to have reflexivity, symmetry, or
transitivity. Define, with a fresh bound label $z$,
\[
 \Sub(x,s):=\forall z\,[\Mem(z,x)\to\Mem(z,s)].
\]
For a property $\phi(x)$, allowing additional parameters, put
\[
 \Upp(x,o):=\Mem(x,o)\to\phi(x),\qquad
 \Low(x;s,o):=\phi(x)\to\bigl(\Sub(x,s)\to\Mem(x,o)\bigr).
\]
The weak-cover assertion at source $s$ is
\begin{equation}\label{eq:weak-cover}
 \exists o\,[\forall x\,\Upp(x,o)\ \wedge\
                 \forall x\,\Low(x;s,o)],
\end{equation}
where $o$ is fresh in the property and source. There are two useful
presentations with an additional free point $p$:
\begin{align}
 W_{\Low}(s,p)&:\quad
 \exists o\,[\forall x\,\Upp(x,o)\wedge
                  \forall x\,\Low(x;s,o)\wedge\Low(p;s,o)],
                  \label{eq:free-lower}\\
 W_{\Upp}(s,p)&:\quad
 \exists o\,[\forall x\,\Upp(x,o)\wedge
                  \forall x\,\Low(x;s,o)\wedge\Upp(p,o)].
                  \label{eq:free-upper}
\end{align}
All replacements of displayed bound letters avoid capture.

\begin{proposition}[Ordinary logical equivalence]\label{prop:free-instance}
In ordinary first-order logic, each of
\eqref{eq:free-lower} and~\eqref{eq:free-upper} is equivalent to
\eqref{eq:weak-cover}.
\end{proposition}
\begin{proof}
For a fixed existential witness $o$, instantiate the corresponding
universal conjunct at $p$ and retain that instance as the third conjunct.
Conversely, discard the third conjunct. No set-theoretic axiom or equality
principle is used.
\end{proof}

In the PA interpretation by binary digits, a witness to
\eqref{eq:weak-cover} is the finite set of all subsets of $s$ satisfying
$\phi$. Each subset, viewed as a bit code, is at most $s$. PA's induction
constructs the bounded bit-coded comprehension for each fixed formula
$\phi$. The extra free conjunct is already true of that witness. This is
an ordinary mathematical model construction; exporting a uniform
arithmetized proof transformation is an additional obligation.

\subsection{The retained numeric rules}

For the bootstrap we retain PA196's six logical/equality heads other than
the choice of orientation of the equality witness, its restricted
implicational basis, modus ponens, and closed generalization. The two
direct membership rewrites are
\begin{align}
 \Eq(v,w)&\to\bigl(\Mem(v,d)\to\Mem(w,d)\bigr),\label{eq:direct-first}\\
 \Eq(v,w)&\to\bigl(\Mem(d,v)\to\Mem(d,w)\bigr).\label{eq:direct-second}
\end{align}
We also retain atomic membership vacuity and add equality vacuity:
\[
 B\to\forall z B
 \quad\text{for }B=\Mem(v,w)\text{ or }\Eq(v,w),
 \qquad z\ne v,w.
\]
These formulas are obtained by projection from guarded bundles; the
numeric guards are the corresponding inequalities of variable labels.

The existing h1 and restricted implicational basis provide explosion and
classical implicational reasoning. The h6 head and generalization lift
proved implications under a universal quantifier. The h2 head exchanges
distinct quantifier dimensions. As in PA196's vacuity proof, these rules
extend atomic vacuity to formulas whose free and bound indices exclude
the chosen fresh label. They also give the following restricted
existential elimination: from $\forall y(A\to C)$ and $\exists y A$,
one may infer $C$ when $C\to\forall y C$ is derivable. None of these
facts assumes a general substitution rule for primitive equality.

The orientation of the retained witness head is essential. Write
\[
 \mathsf{W_R}:\quad\exists y\,\Eq(v,y),
 \qquad
 \mathsf{W_L}:\quad\exists y\,\Eq(y,v).
\]
Here $\mathsf{W_R}$ is PA196's mirrored witness; $\mathsf{W_L}$ is the
original orientation. The free-lower construction below uses
$\mathsf{W_R}$, and the free-upper construction uses $\mathsf{W_L}$.
The other pairing of a free instance with a witness orientation is not
certified by the argument in this section.

\subsection{Two restricted specialization lemmas}

Let $\chi(t)$ be one of
$\Mem(d,t)$, $\Mem(t,d)$, or $\Mem(t,t)$. The two direct rewrites give
\[
 \Eq(v,w)\to\bigl(\chi(v)\to\chi(w)\bigr).
\]
For the diagonal atom, apply the two coordinate rewrites in succession.
All freshness conditions stated below concern numeric variable labels.

\begin{lemma}[Negative upper specialization]\label{lem:negative-upper}
With $\mathsf{W_R}$, if $y\ne v,o$ and $d\ne y$ when $d$ occurs, then
\[
 \forall y\,[\Mem(y,o)\to\neg\chi(y)]
 \ \to\ [\Mem(v,o)\to\neg\chi(v)]
\]
is derivable.
\end{lemma}
\begin{proof}
Direct rewriting gives, under $\Eq(v,y)$, both
$\Mem(v,o)\to\Mem(y,o)$ and $\chi(v)\to\chi(y)$. Classical reasoning
therefore gives
\[
 \Eq(v,y)\to
 \bigl((\Mem(y,o)\to\neg\chi(y))\to
       (\Mem(v,o)\to\neg\chi(v))\bigr).
\]
Generalize at $y$ and use h6 with the displayed universal premise.
The conclusion $\Mem(v,o)\to\neg\chi(v)$ is vacuous at $y$.
Eliminate the witness $\exists y\,\Eq(v,y)$ by the restricted
existential rule. This is a proof of this particular specialization,
not an invocation of unrestricted universal instantiation.
\end{proof}

\begin{lemma}[Negative lower specialization]\label{lem:negative-lower}
With $\mathsf{W_L}$, let $\phi(t)=\neg\chi(t)$. If $y$ differs from
$v,s,o$ and from $d$ when present, and the subset binder $o$ is fresh for
$y,v,s$, then
\[
 \forall y\,\Low(y;s,o)\ \to\ \Low(v;s,o)
\]
is derivable. The same letter $o$ may name the free output in the
membership atom and the locally bound variable inside $\Sub$.
\end{lemma}
\begin{proof}
Under $\Eq(y,v)$, the direct rewrites imply
\[
 \neg\chi(v)\to\neg\chi(y),\qquad
 \Mem(y,o)\to\Mem(v,o).
\]
They also imply $\Sub(v,s)\to\Sub(y,s)$ under that equality
assumption. Indeed,
$\Eq(y,v)\to(\Mem(o,y)\to\Mem(o,v))$ is a direct rewrite.
Generalization at $o$, equality vacuity at $o$, and h6 lift it to
$\Eq(y,v)\to\Sub(y,v)$. Combine it pointwise with
$\Sub(v,s)$, using h6, to obtain $\Sub(y,s)$. Consequently
\[
 \Eq(y,v)\to\bigl(\Low(y;s,o)\to\Low(v;s,o)\bigr)
\]
is derivable. Generalize at $y$, combine with $\forall y\Low(y;s,o)$,
and eliminate $\exists y\Eq(y,v)$ using $\mathsf{W_L}$ and vacuity of
$\Low(v;s,o)$ at $y$. Only the explicitly displayed rewrites and
quantifier rules have been used.
\end{proof}

\subsection{Recovery of reverse membership rewriting}

\begin{proposition}[Matching free instances repair the bootstrap]
\label{prop:reverse-repair}
Either of the following additions to the retained rules derives both
reverse membership rewrites:
\begin{enumerate}
 \item $W_{\Low}$ together with $\mathsf{W_R}$;
 \item $W_{\Upp}$ together with $\mathsf{W_L}$.
\end{enumerate}
More precisely, for every $a,b,d$ they derive
\[
 \Eq(a,b)\to(\Mem(d,b)\to\Mem(d,a)),\qquad
 \Eq(a,b)\to(\Mem(b,d)\to\Mem(a,d)).
\]
The old primitive equality-transitivity head is not used.
\end{proposition}
\begin{proof}
For $a=b$ the conclusions are propositional tautologies. Suppose
$a\ne b$.

First use $W_{\Low}$ and $\mathsf{W_R}$. For $d\ne a$, take
$\phi(t)=\neg\chi(t)$, source $s=b$, and distinguished point $p=a$.
Choose the output label $o$ fresh for $a,b,d$.
By direct rewriting, equality vacuity, generalization and h6,
\[
 \Eq(a,b)\to\Sub(a,b).
\]
Assume $\Eq(a,b)$, $\chi(b)$ and $\neg\chi(a)$. The free lower
conjunct gives $\Mem(a,o)$, the direct first-coordinate rewrite gives
$\Mem(b,o)$, and Lemma~\ref{lem:negative-upper}, applied to the upper
universal conjunct, gives $\neg\chi(b)$. This is a contradiction.
The final implication has no index $o$, so its vacuity permits
elimination of the existential output. Thus
$\Eq(a,b)\to(\chi(b)\to\chi(a))$ is obtained for $d\ne a$ and for
the diagonal atom.
If $d=a$, each of $\Mem(a,b)$ and $\Mem(b,a)$ transfers directly to
$\Mem(b,b)$ under $\Eq(a,b)$. The diagonal reverse implication gives
$\Mem(a,a)$, resolving both omitted cases.

Next use $W_{\Upp}$ and $\mathsf{W_L}$. For $d\ne b$, take
$\phi(t)=\neg\chi(t)$, source $s=a$, and distinguished point $p=b$,
again with a fresh output $o$. Under $\Eq(a,b)$, $\chi(b)$ and
$\neg\chi(a)$, Lemma~\ref{lem:negative-lower} specializes the lower
universal conjunct to $a$. The formula $\Sub(a,a)$ follows by
generalizing an implicational identity. Hence $\Mem(a,o)$ follows.
Direct rewriting gives $\Mem(b,o)$, and the free upper conjunct at
$b$ gives $\neg\chi(b)$, a contradiction. Eliminate the existential
output as before. For $d=b$, the diagonal reverse implication takes
$\Mem(b,b)$ to $\Mem(a,a)$; the direct rewrites then yield
$\Mem(b,a)$ and $\Mem(a,b)$ respectively, settling the two omitted
cases.
\end{proof}

This provides the two directions of membership rewriting needed for
PA196's defined-neighborhood equality. Generalization and equality
vacuity then give the implication from primitive $\Eq$ to that defined
relation. The relevant equality witness is transported to the defined
relation by existential monotonicity, reversing the \emph{defined}
relation when necessary. No claim that primitive $\Eq$ itself has a
separate symmetry or transitivity theorem is required.

\subsection{Native codes and the remaining obligations}

The argument respects the special coupling between a property code and
an output-variable label. In the U208 role-emitter notation, the same
number $a$ is both the property code and the output label; $b$ is the
source label and $c$ the tested label. Put
\begin{align*}
 e&=\Mem(c,a), &
 G&=\forall a\,[\Mem(a,c)\to\Mem(a,b)],\\
 S&=e\to a, & U&=\forall c\,S,\\
 R&=a\to(G\to e), & L&=\forall c\,R.
\end{align*}
Here $a$ in formula position denotes the decoded formula with code $a$.
The two revised cover tails are
\[
 H_R=\forall a\,[U\to(L\to(R\to\bot))],\qquad
 H_S=\forall a\,[U\to(L\to(S\to\bot))].
\]
With $\Delta=\Eq(a,b)\to\forall c\Eq(a,b)$ and
$V=\Mem(a,b)\to\forall c\Mem(a,b)$, emit respectively
\[
 \neg\bigl(\Delta\to(V\to H_R)\bigr),\qquad
 \neg\bigl(\Delta\to(V\to H_S)\bigr)
\]
under the original guard $c\ne a,b$. Keep truth as the disabled output.
The free-lower version uses the mirrored equality witness, and the
free-upper version uses the original witness.

For the negative-property instances used in the proof, the property code
exceeds its displayed atom indices. If it equals the other distinguished
label, use the equivalent triple negation instead of a single negation;
its code is strictly larger. The contraposition argument still supplies
the same negative-property specialization, and the output label is now
fresh for all indices in the final conclusion. The source and test labels
are different because the case $a=b$ in the proof was already discharged
propositionally. Thus the native guard does not remove any case needed
for Proposition~\ref{prop:reverse-repair}.

In ordinary first-order logic the new tails add only an instance of an
existing universal conjunct. This observation also covers the original
role-emitter's exceptional equality of its source and output labels: the
native free instance has the same local binding as the body it copies.
The separate calibration of that exceptional old weak-cover instance
must still be retained.

\paragraph{Scope.}
The new component is a noncircular derivation of the reverse-rewriting
interface for two precisely specified emitter variants. A complete
construction must still supply the exact AX stack contract and selection
ranks; complete numeric formula/proof translations, including any
argument or tag recoding; a physical allocation with correct startup and
universal runtime contracts; and both coded trace/proof directions in
PA. This section supplies none of those by implication. In particular,
it establishes neither a machine with fewer than 196 states nor the full
PA-provable nonhalting--consistency biconditional.

\endgroup
